\begin{abstract}
Deployed agents run under operator-assigned personas, alongside other agents,
and with whatever tools they are given. We ask whether these change whether an
agent follows a stated rule. In \R{all}{nscenarios} scenarios, an operations
agent can finish its task only by sending a network request its policy
forbids. We cross the tool it is given (a named HTTP tool or a code sandbox)
with four operator personas and four kinds of peer-agent testimony, across
\R{all}{nmodels} models (\R{all}{nsamples} episodes). Every measure is computed
mechanically from the transcript, without an LLM judge, and a permit-policy
control separates compliance from incapacity. Four frontier models almost
never broke the rule under any condition, although they made the request
whenever it was permitted. Among the \R{all}{nactive} models that did break
it, the same act was far more likely when written as code than as a named tool
call (pooled OR \R{pooled}{Honeor}); peer notes reporting a violation raised
the violation rate (OR \R{pooled}{Htwoaor}) and notes reporting a refusal
lowered it (OR \R{pooled}{Htwobor}); and a ``cautious'' persona lowered it (OR
\R{pooled}{Hthreebor}) while a ``resourceful'' one raised it (OR
\R{pooled}{Hthreeaor}). These pooled effects hide differences in direction
between models: one model reverses both peer effects, and the same persona
raises violation in one model and lowers it in another. Rule-following
evaluations should therefore vary tool and social context, and report effects
per model.
\end{abstract}
